\chapter{Capability Monitoring Specification}\label{ch:monitoring}

This chapter specifies the emergent capability monitoring system, addressing the
IMDA MGF requirement (Section 2.3.3) that ``monitoring for emergent capabilities
must be implemented with defined detection triggers.''

\section{Overview}

AI systems may develop unexpected capabilities or behaviours through:
\begin{itemize}
  \item Model updates that alter response patterns
  \item Prompt injection or adversarial inputs
  \item Emergent behaviours from complex interactions
  \item Drift in input data distributions
\end{itemize}

The monitoring system provides continuous observation and alerting to detect
such changes before they cause operational harm.

\section{Architecture}

\begin{figure}[htbp]
\centering
\begin{tikzpicture}[
  node distance=10mm and 18mm,
  box/.style={draw,rounded corners,minimum width=32mm,minimum height=12mm,align=center,font=\small,fill=blue!5},
  mon/.style={draw,rounded corners,minimum width=32mm,minimum height=12mm,align=center,font=\small,fill=orange!10},
  alert/.style={draw,rounded corners,minimum width=28mm,minimum height=10mm,align=center,font=\small,fill=red!10},
  db/.style={draw,cylinder,shape border rotate=90,aspect=0.25,minimum width=20mm,minimum height=14mm,fill=gray!10,font=\small},
  t/.style={-{Latex[length=2mm]}}
]
\node[box] (pipeline) {AI Pipeline};
\node[mon,right=of pipeline,align=center] (collector) {Metrics\\Collector};
\node[mon,right=of collector,align=center] (analyzer) {Anomaly\\Analyzer};
\node[alert,below=of analyzer,align=center] (alert) {Alert\\Manager};
\node[db,below=of collector,align=center] (store) {Metrics\\Store};

\draw[t] (pipeline) -- node[above,font=\scriptsize]{telemetry} (collector);
\draw[t] (collector) -- (analyzer);
\draw[t] (collector) -- (store);
\draw[t] (analyzer) -- (alert);
\draw[t] (store.east) -- (analyzer.south west);
\end{tikzpicture}
\caption{Capability monitoring architecture.}
\label{fig:monitor-arch}
\end{figure}

\section{Metrics collection}

\subsection{Required metrics}

\begin{table}[htbp]
\centering
\caption{Capability monitoring metrics.}
\label{tab:monitor-metrics}
\footnotesize
\begin{tabularx}{\linewidth}{@{}l l X@{}}
\toprule
\textbf{Metric} & \textbf{Type} & \textbf{Description} \\ \midrule
\texttt{output\_token\_count} & histogram & Distribution of output lengths \\
\texttt{response\_latency\_ms} & histogram & Response time distribution \\
\texttt{confidence\_score} & histogram & AI confidence distribution \\
\texttt{refusal\_rate} & counter & Rate of model refusals \\
\texttt{tool\_call\_frequency} & counter & Frequency of tool/function calls \\
\texttt{output\_format\_violations} & counter & JSON schema validation failures \\
\texttt{hallucination\_flags} & counter & Detected hallucination instances \\
\texttt{unexpected\_topic\_detections} & counter & Off-topic responses \\
\texttt{novel\_capability\_indicators} & gauge & Composite novelty score \\
\bottomrule
\end{tabularx}
\end{table}

\subsection{Collection points}

Metrics are collected at the following pipeline stages:

\begin{enumerate}
  \item \textbf{Pre-inference}: Input characteristics (length, complexity, topic)
  \item \textbf{Post-inference}: Output characteristics (length, confidence, format)
  \item \textbf{Post-validation}: Validation results (schema, business rules)
  \item \textbf{Post-review}: Human review outcomes (accept/edit/reject rates)
\end{enumerate}

\section{Detection triggers}

\subsection{Statistical anomaly detection}

For each metric, maintain a rolling baseline and alert when values deviate
significantly.

\begin{table}[htbp]
\centering
\caption{Anomaly detection thresholds.}
\label{tab:anomaly-thresholds}
\footnotesize
\begin{tabularx}{\linewidth}{@{}l l X@{}}
\toprule
\textbf{Metric} & \textbf{Alert Threshold} & \textbf{Baseline Period} \\ \midrule
Output length (mean) & $>2\sigma$ shift & 7-day rolling \\
Latency (p95) & $>50\%$ increase & 24-hour rolling \\
Confidence (mean) & $>10\%$ drop & 7-day rolling \\
Refusal rate & $>3\times$ baseline & 7-day rolling \\
Schema violations & $>5\%$ of requests & 24-hour rolling \\
\bottomrule
\end{tabularx}
\end{table}

\subsection{Behavioural pattern detection}

Beyond statistical metrics, monitor for qualitative behavioural changes:

\begin{enumerate}
  \item \textbf{Topic drift}: Cosine similarity of output embeddings to expected
        topics falls below 0.7
  \item \textbf{Instruction leakage}: Presence of system prompt fragments in output
  \item \textbf{Novel tool usage}: Model attempts to call undeclared tools
  \item \textbf{Self-referential loops}: Model references its own previous outputs
        in unexpected ways
\end{enumerate}

\subsection{Trigger definitions}

\begin{lstlisting}[language=Python,basicstyle=\ttfamily\small,caption={Detection trigger configuration.}]
triggers = {
    "output_length_anomaly": {
        "metric": "output_token_count",
        "condition": "mean > baseline_mean + 2 * baseline_std",
        "severity": "WARNING",
        "cooldown_minutes": 60
    },
    "latency_spike": {
        "metric": "response_latency_ms",
        "condition": "p95 > 1.5 * baseline_p95",
        "severity": "WARNING",
        "cooldown_minutes": 15
    },
    "high_refusal_rate": {
        "metric": "refusal_rate",
        "condition": "rate > 3 * baseline_rate",
        "severity": "CRITICAL",
        "cooldown_minutes": 5
    },
    "schema_violation_surge": {
        "metric": "output_format_violations",
        "condition": "rate > 0.05",  # >5% failure
        "severity": "CRITICAL",
        "cooldown_minutes": 5
    },
    "topic_drift": {
        "metric": "topic_similarity",
        "condition": "mean < 0.7",
        "severity": "WARNING",
        "cooldown_minutes": 30
    }
}
\end{lstlisting}

\section{Alert management}

\subsection{Alert severity levels}

\begin{table}[htbp]
\centering
\caption{Alert severity definitions.}
\label{tab:alert-severity}
\footnotesize
\begin{tabularx}{\linewidth}{@{}l l X@{}}
\toprule
\textbf{Severity} & \textbf{Response Time} & \textbf{Action} \\ \midrule
CRITICAL & Immediate & Page on-call; consider circuit breaker \\
WARNING & 1 hour & Investigate during business hours \\
INFO & Next business day & Review in weekly monitoring meeting \\
\bottomrule
\end{tabularx}
\end{table}

\subsection{Alert routing}

\begin{lstlisting}[basicstyle=\ttfamily\small,caption={Alert routing configuration.}]
routing:
  CRITICAL:
    channels:
      - pagerduty: "ai-ops-critical"
      - slack: "#ai-incidents"
      - email: "ai-ops@company.com"
    escalation_after_minutes: 15
    
  WARNING:
    channels:
      - slack: "#ai-monitoring"
      - email: "ai-ops@company.com"
    escalation_after_minutes: 60
    
  INFO:
    channels:
      - slack: "#ai-monitoring"
    escalation_after_minutes: null  # no escalation
\end{lstlisting}

\subsection{Circuit breaker integration}

When CRITICAL alerts persist, the monitoring system can activate a circuit
breaker:

\begin{enumerate}
  \item After 3 consecutive CRITICAL alerts within 15 minutes
  \item Circuit breaker opens: all AI requests routed to manual queue
  \item Alert sent to AI governance team
  \item Manual intervention required to close circuit breaker
\end{enumerate}

\subsection{Graceful degradation ladder}
\label{sec:degradation-ladder}

A hard circuit-break (manual queue only) is the last resort. Before
that point, the monitoring system walks a four-rung degradation
ladder, chosen so that each rung trades off quality for availability
in a \emph{controlled and observable} way.
Table~\ref{tab:degradation-ladder} specifies the ladder per
pipeline; rung escalation is automatic, rung de-escalation requires
a human ack.

\begin{table}[htbp]
\centering
\caption{Graceful degradation ladder (from normal operation down
to circuit break). Transitions are logged as governance events.}
\label{tab:degradation-ladder}
\footnotesize
\begin{tabularx}{\linewidth}{@{}l l l X@{}}
\toprule
\textbf{Rung} & \textbf{Name} & \textbf{Trigger} & \textbf{Action} \\ \midrule
0 & Normal & p95 latency $<$ SLO, error rate $<$ 1\,\% & Primary model on primary GPU (e.g.\ Qwen2.5-32B on A100). \\
1 & Reduced concurrency & p95 latency $\ge$ SLO for 5\,min & Halve request concurrency; queue depth bounded; user-visible latency warning. \\
2 & Smaller model & Primary model unavailable or OOM for 2\,min & Automatic failover to smaller model (TB: 32B$\to$14B; Arabic AP: 14B$\to$7B; Simula: 32B$\to$14B for non-user-facing batches). Output tagged \texttt{model\_tier=fallback} for audit. \\
3 & Cached / heuristic & Smaller model also unhealthy & Return cached response if cache hit within TTL; else return heuristic output (TB: statistical threshold only; Arabic AP: OCR-only, no translation; Simula: block generation). \\
4 & Circuit break & Rung~3 still failing for 5\,min & All AI calls routed to manual queue; page on-call; governance ticket. \\
\bottomrule
\end{tabularx}
\end{table}

\paragraph{Tagging and auditability.}
Every response served from rung~1--3 carries a \texttt{model\_tier}
attribute in its envelope and an audit event referencing the rung
and the ladder-transition decision. Downstream consumers (review
dashboard, AP workflow) render a visible ``fallback'' badge so
reviewers know the response did not come from the primary model.
This preserves the IMDA MGF~2.1.2-002 attribution guarantee across
all rungs: attribution is to the \emph{model actually used}, not
to the model the request was routed to.

\paragraph{Rung selection rationale.}
Rung~1 (concurrency cut) is cheap and preserves quality; it buys
time for transient load spikes. Rung~2 (smaller model) is the
primary availability lever because model-host failures are the
most common outage mode. Rung~3 (cached/heuristic) is the last
value-preserving rung. Rung~4 is reserved for the case where the
AI subsystem cannot be trusted at all; it is deliberately
\emph{not} automatic to close --- human judgement is required to
confirm the root cause is resolved. Rung transitions upward are
automatic; transitions downward (closer to normal) require either
a human acknowledgement or a clean 15-minute signal window.

\section{Dashboard specification}

\subsection{Real-time view}

The monitoring dashboard provides:

\begin{itemize}
  \item Live metric charts (1-minute granularity)
  \item Active alert panel with acknowledge/resolve actions
  \item System health summary (green/yellow/red)
  \item Recent trigger history
\end{itemize}

\subsection{Historical analysis}

For trend analysis and post-incident review:

\begin{itemize}
  \item Metric time series with configurable ranges
  \item Alert history with resolution notes
  \item Correlation views (multiple metrics overlaid)
  \item Baseline comparison charts
\end{itemize}

\section{Audit and compliance}

\subsection{Retention requirements}

\begin{table}[htbp]
\centering
\caption{Monitoring data retention.}
\label{tab:monitor-retention}
\footnotesize
\begin{tabularx}{\linewidth}{@{}l r X@{}}
\toprule
\textbf{Data Type} & \textbf{Retention} & \textbf{Storage} \\ \midrule
Raw metrics & 90 days & Time-series database \\
Aggregated metrics & 2 years & Data warehouse \\
Alert history & 7 years & Compliance archive \\
Incident reports & 7 years & Document management \\
\bottomrule
\end{tabularx}
\end{table}

\subsection{Regulatory reporting}

For IMDA MGF compliance, generate quarterly reports containing:

\begin{enumerate}
  \item Summary of monitoring coverage (metrics, triggers, alerts)
  \item List of all CRITICAL alerts and resolutions
  \item Trend analysis of key capability metrics
  \item Any identified emergent behaviours and mitigations
  \item Circuit breaker activations and durations
\end{enumerate}

\section{Implementation checklist}

\begin{itemize}
  \item[\texttt{[ ]}] Metrics collector integrated with all AI pipelines
  \item[\texttt{[ ]}] All metrics in Table~\ref{tab:monitor-metrics} collected
  \item[\texttt{[ ]}] Anomaly detection for all triggers in Table~\ref{tab:anomaly-thresholds}
  \item[\texttt{[ ]}] Alert routing configured for all severity levels
  \item[\texttt{[ ]}] Circuit breaker integration tested
  \item[\texttt{[ ]}] Dashboard with real-time and historical views
  \item[\texttt{[ ]}] Retention policies implemented per Table~\ref{tab:monitor-retention}
  \item[\texttt{[ ]}] Quarterly report template created
  \item[\texttt{[ ]}] Runbook for each alert type documented
  \item[\texttt{[ ]}] Integration tests with simulated anomalies
\end{itemize}
